\section{Set Theory}
\label{sec:set-theory}

\emph{Zermelo--Fraenkel set theory with the axiom of choice} (ZFC) is the
standard first-order foundation of modern mathematics, and the theory in which
the developments described in this thesis are ultimately expressed. We
recall it here to fix the axioms and the vocabulary used in later chapters.
After sketching how ZFC emerged from Zermelo's original axiomatization to its
modern Zermelo--Fraenkel form, we state the theory formally and develop the
derived notions of ordered pairs, relations, and functions. This section is a
condensed introduction; Jech~\cite{jech2003set} and Kunen~\cite{kunen2011set}
are standard, accessible textbooks on the subject.

\subsection{From Naive to Axiomatic Set Theory}
\label{subsec:set-theory-history}

Set theory originates with Georg Cantor, who in the 1870s and 1880s was led, by
his study of trigonometric series, to treat infinite collections as completed
mathematical objects and to develop a theory of their sizes through cardinal and
ordinal numbers. Cantor's informal conception is captured by his well-known
description of a set as ``any collection into a whole of definite, distinct
objects of our intuition or thought''~\cite{CSB_thm}. Two principles
are implicit in this view: \emph{extensionality}, that a set is wholly
determined by its members; and \emph{unrestricted comprehension}, that every
property~$\varphi$ delimits a set $\{x \mid \varphi(x)\}$ of all objects
satisfying it. This naive theory was extraordinarily productive and, taken at
face value, inconsistent.

The inconsistency was made explicit by Russell in 1901~\cite{russell1903principles}.
Applying comprehension to the property $x \notin x$ yields the
\emph{Russell set} $E = \{x \mid x \notin x\}$; asking whether $E \in E$ gives
\[
  E \in E \iff E \notin E,
\]
a contradiction, today commonly called \emph{Russell's paradox}.
Related antinomies---Cantor's paradox of the largest cardinal
and the Burali-Forti paradox of the largest ordinal---showed that the difficulty
was not an isolated curiosity but a structural defect of unrestricted
comprehension~\cite{vanHeijenoort1967}. A collection cannot be formed merely
because a property describes it.

The lasting response was to replace comprehension with axioms that delimit
exactly which collections may be formed, while retaining enough strength to
reconstruct ordinary mathematics. This is the axiomatic route opened by Zermelo
and completed by Fraenkel, Skolem, and von Neumann, which we now follow.

\subsection{Zermelo's Axiomatization}
\label{subsec:set-theory-zermelo}

The first successful axiomatization is due to
Zermelo~\cite{zermelo1908untersuchungen}. It was motivated in part by the need to
secure his 1904 proof of the well-ordering theorem and the axiom of choice on
which it relied, then a matter of controversy. Zermelo's system comprises seven
axioms: \emph{extensionality}; \emph{elementary sets}, providing the empty set
and unordered pairs; \emph{separation}; \emph{power set}; \emph{union};
\emph{choice}; and \emph{infinity}.

Its decisive idea is the axiom of \emph{separation} (\emph{Aussonderung}), which
tames comprehension by relativizing it to a set already given: from a set $x$
and a property $\varphi$ one may form the subset
$\{z \in x \mid \varphi(z)\}$, but one may not conjure
$\{z \mid \varphi(z)\}$ out of nothing. Russell's paradox thereby dissolves. The
collection $\{z \in x \mid z \notin z\}$ is a harmless set, and the diagonal
argument that produced $E$ now establishes only that this subset does not belong
to $x$; hence no set $x$ contains all sets, and the universe of sets is not
itself a set.

Zermelo's theory has two well-known shortcomings, both addressed within the
following two decades. First, the notion of a ``definite'' property~$\varphi$
admissible in separation was left informal, with no precise delimitation of the
permitted properties. Second, the system is too weak: it cannot prove the
existence of the set $\{Z_0, Z_1, Z_2, \dots\}$, where $Z_0 = \mathbb{N}$ and
$Z_{n+1} = \pow(Z_n)$, nor of cardinals such as $\aleph_\omega$. Sets whose
construction requires iterating an operation along an arbitrary index lie beyond
its reach.

\subsection{Replacement, Foundation, and the Zermelo--Fraenkel System}
\label{subsec:set-theory-zf}

Both shortcomings were remedied, independently, by
Fraenkel~\cite{fraenkel1922grundlagen} and Skolem~\cite{skolem1922bemerkungen}
around 1922. Skolem proposed making ``definite property'' precise as a
first-order formula in the language with equality and membership, thereby fixing
the meaning of separation. They both added the axiom schema of
\emph{replacement}: if a formula~$\varphi$ defines a functional relation, then
the image of any set under it is again a set. Replacement supplies exactly the
missing closure---$\{Z_n \mid n \in \mathbb{N}\}$ is the image of $\mathbb{N}$
under $n \mapsto Z_n$---and with it the transfinite recursion and ordinal
arithmetic that Zermelo's system could not justify.

A final axiom, \emph{foundation} (or \emph{regularity}), was isolated by von
Neumann~\cite{vonNeumann1925axiomatisierung} and incorporated into the system by
Zermelo~\cite{zermelo1930grenzzahlen} in 1930. It asserts that membership is
well-founded, outlawing infinite descending chains
$\cdots \in x_2 \in x_1 \in x_0$ and, in particular, sets satisfying $x \in x$.
Foundation introduces no new sets; instead it disciplines the universe,
stratifying every set into the \emph{cumulative hierarchy} obtained by iterating
the power-set operation through the ordinals (\Cref{subsec:set-theory-hierarchy}).
This well-founded picture is precisely the one that the Lean model of
\Cref{ch:zflean} reflects, where sets are realized as a well-founded inductive
type.

The resulting first-order theory is \emph{Zermelo--Fraenkel set theory} (ZF);
adjoining the axiom of choice yields ZFC. It is formulated in classical
first-order logic with equality over a single binary relation symbol~$\in$, and
has become the de facto foundation of mainstream mathematics. We state it
precisely.

\subsection{The Axioms of ZF}
\label{subsec:set-theory-axioms}

We work in classical first-order logic with equality, whose only non-logical
symbol is the binary relation~$\in$; every object of the theory is a \emph{set},
so the elements of a set are themselves sets and the universe contains nothing
else (\emph{pure} set theory).

We state the axioms in ordinary mathematical style rather than fully formal
syntax, assuming the usual high-level handling of variables: bound variables are
taken suitably fresh, any free parameters of a formula are read as implicitly
universally quantified, and we use standard abbreviations throughout.
In particular, $x \regnotation{\subseteq} y$ stands for $\forall z,\, z \in x \to z \in y$
and $x \neq y$ for $\lnot(x = y)$; the bounded quantifiers $\forall z \in x,\,\psi$
and $\exists z \in x,\,\psi$ abbreviate $\forall z,\,(z \in x \to \psi)$ and
$\exists z,\,(z \in x \land \psi)$; and $\regnotation{\varnothing}$, $\regnotation{\{y\}}$,
$\regnotation{x \cup y}$, $\regnotation{x \cap y}$ denote the empty set, singleton,
binary union, and binary intersection, each unfolding into a formula over $\in$
alone.\footnote{%
  Strictly, the language of ZF has no function symbols, only the relation
  $\in$, so these notations are not genuine terms but \emph{eliminable
  abbreviations}, each justified by a provable unique-existence statement
  (\eg $\exists! w,\,\forall z,\,z \in w \leftrightarrow (z = x \lor z = y)$ for
  $\{x, y\}$): a formula mentioning one stands for the $\in$-formula obtained by
  unfolding it in context. We adopt this reading throughout, rather than treating
  the notations as primitive terms via a definite-description operator.}
The axioms of ZF are the following eight; the ninth, choice, stands apart
and is presented in \Cref{subsec:set-theory-choice}.

\begin{enumerate}[label=\textbf{(A\arabic*)}, ref=A\arabic*, leftmargin=3.6em, labelsep=1em, itemsep=0.8ex]

  \item\label{ax:ext} \axiomlayout{%
    \textbf{Extensionality.} Sets with the same elements are equal:
    \axiomeq{\forall x\, y,\,(\forall z,\,z \in x \leftrightarrow z \in y) \to x = y}
  }{\axfigExt}

  \item\label{ax:pair} \axiomlayout{%
    \textbf{Pairing.} For all $x$ and $y$ there is a set whose elements are exactly $x$ and $y$:
    \axiomeq{\forall x\, y,\,\exists z,\,\forall w,\,w \in z \leftrightarrow (w = x \lor w = y)}
  }{\axfigPair}
  \begin{note}
    Such a set \(z\) is unique by extensionality, and is denoted \(\regnotation{\{x, y\}}\).
  \end{note}

  \item\label{ax:union} \axiomlayout{%
    \textbf{Union.} For every set $x$ there is a set \(y\) whose elements are the elements of the elements of $x$:
    \axiomeq{\forall x,\,\exists y,\,\forall z,\,z \in y \leftrightarrow (\exists w,\,w \in x \land z \in w)}
  }{\axfigUnion}
  \begin{note}
    Such a set \(y\) is unique by extensionality, and is denoted \(\regnotation{\bigcup} x\).
  \end{note}

  \item\label{ax:pow} \axiomlayout{%
    \textbf{Power set.} For every set $x$ there is a set \(y\) whose elements are exactly the subsets of $x$:
    \axiomeq{\forall x,\,\exists y,\,\forall z,\,z \in y \leftrightarrow z \subseteq x}
  }{\axfigPow}
  \begin{note}
    Such a set \(y\) is unique by extensionality, and is denoted \(\regnotation{\mathcal{P}}(x)\).
  \end{note}

  \item\label{ax:inf} \axiomlayout{%
    \textbf{Infinity.} There exists an \emph{inductive} set, \ie one containing $\varnothing$ and closed under the successor $y \mapsto y \cup \{y\}$:
    \axiomeq{\exists x,\,\varnothing \in x \;\land\; (\forall y,\,y \in x \to y \cup \{y\} \in x)}
  }{\axfigInf}

  \item\label{ax:found} \axiomlayout{%
    \textbf{Foundation.} Every non-empty set has a member disjoint from it; equivalently, $\in$ is well-founded:
    \axiomeq{\forall x,\, x \neq \varnothing \;\to\; (\exists y,\,y \in x \land y \cap x = \varnothing)}
  }{\axfigReg}

  \item\label{ax:sep} \axiomlayout{%
    \textbf{Separation} \textup{(schema).} For every formula $\varphi(z)$, the subset $y$ of $x$ cut out by $\varphi$ exists:
    \axiomeq{\forall x,\,\exists y,\,\forall z,\;z \in y \leftrightarrow (z \in x \land \varphi(z))}
  }{\axfigSep}
  \begin{note}
    Such a set \(y\) is unique by extensionality, and is denoted \(\regnotation{\setcomp{z \in x}{\varphi(z)}}\).
  \end{note}

  \item\label{ax:repl} \axiomlayout{%
    \textbf{Replacement} \textup{(schema).} For any definable\footnotemark{} operation $x \mapsto f(x)$, the image of any set under it is a set:
    \axiomeq{\forall a,\,\exists b,\,\forall y,\; (y \in b \leftrightarrow \exists x,\,x \in a \land f(x) = y)}
  }{\axfigRepl}
  \footnotetext{\emph{Definable} means given by a first-order formula $f(x, y)$ over~$\in$ (possibly with set parameters) relating each $x$ to at most one $y$; replacement is thus a schema, with one instance per such formula.}

  \begin{note}
    The resulting set \(b\) is the \emph{image} of \(a\) under \(f\), denoted \(\regnotation{\setcomp{f(x)}{x \in a}}\).
  \end{note}
\end{enumerate}

\paragraph{Derived constructions.}
A few standard constructions are worth mentioning, since they are part of the
common vocabulary of set theory.
Dually to union, the \emph{intersection} of a non-empty set $x$ defined as a subset
of its union cut out by the property of belonging to every member of $x$, using
union \eqref{ax:union} and separation \eqref{ax:sep}:
\[
\regnotation{\bigcap} x \triangleq \setcomp{z \in \bigcup x}{\forall w \in x,\, z \in w}.
\]
%
Using pairing \eqref{ax:pair}, one can form the \emph{singleton} \(\{x\} \triangleq \{x, x\}\)
as the set whose only element is \(x\).
The \emph{ordered pair}, due to Kuratowski, of $x$ and $y$ is then encoded as the set
$\regnotation{\opair{x}{y}} \triangleq \{\{x\}, \{x, y\}\}$, whose only essential
property is that $\opair{x}{y} = \opair{x'}{y'}$ exactly when $x = x'$ and
$y = y'$.
%
As $\opair{x}{y} \in \pow(\pow(x \cup y))$, separation \eqref{ax:sep} then carves the
\emph{Cartesian product} $A \regnotation{\times} B$ out of $\pow(\pow(A \cup B))$,
\[
  A \times B \triangleq \setcomp{z \in \pow(\pow(A \cup B))}{\exists x \in A\,\exists y \in B,\; z = \opair{x}{y}},
\]
the set of all ordered pairs with first component in $A$ and second in $B$, or equivalently
using \eqref{ax:repl}, the set of all ordered pairs of elements of $A$ and $B$,
\[
  A \times B \triangleq \setcomp{\opair{x}{y}}{x \in A,\,y \in B}.
\]

\paragraph{Redundancies and schemas.}
Several axioms are stated for readability rather than independence. Separation is
a consequence of replacement, and pairing follows from power set and replacement.
They are nonetheless retained as primitive axioms for two reasons: they are
elementary set-formation principles, whose derivation from replacement is an
artificial detour; and they are exactly what remains in the weaker subsystems
obtained by dropping or restricting replacement---notably Zermelo set theory,
which already suffices for large parts of ordinary mathematics and is studied in
its own right. A minimal presentation of ZF therefore retains only
extensionality, union, power set, infinity, foundation, and replacement. Because
first-order logic admits no quantification over formulas, separation and
replacement cannot be single sentences: each is an \emph{axiom schema} denoting
infinitely many instances, one per formula~$\varphi$. This is an essential
feature of the first-order formulation---and the reason ZF, and a fortiori ZFC,
is not finitely axiomatizable.

\subsection{The Axiom of Choice}
\label{subsec:set-theory-choice}

Unlike the existence axioms of ZF, each of which specifies the elements of the
set it asserts to exist, the axiom of choice asserts the existence of a set
without determining its members. Adjoining it to ZF yields \emph{Zermelo--
Fraenkel set theory with choice}, ZFC.

\begin{enumerate}[label=\textbf{(A\arabic*)}, ref=A\arabic*, start=9, leftmargin=3.6em, labelsep=1em, itemsep=0.8ex]

  \item\label{ax:choice} \axiomlayout{%
    \textbf{Choice.} Every set of pairwise-disjoint non-empty sets has a \emph{selector}---a set meeting each of its members in exactly one element:
    \axiomeq{\forall x,\,\mathrm{Disj}(x) \;\to\; \exists c,\,\forall y \in x,\,\exists! z,\,z \in c \land z \in y}
    where $\mathrm{Disj}(x)$ abbreviates
    \[
      \varnothing \notin x \;\land\; \forall y \in x,\,\forall y' \in x,\,(y \neq y' \to y \cap y' = \varnothing).
    \]
  }{\axfigChoice}

\end{enumerate}

The axiom genuinely extends bare ZF, being \emph{independent} of the other eight:
if ZF is consistent, then so is ZF with choice, by G\"odel's constructible
universe~\cite{godel1940consistency}, and so is ZF with the negation of choice, by
Cohen's forcing~\cite{cohen1966settheory}. Choice is therefore neither provable
nor refutable from \eqref{ax:ext}--\eqref{ax:repl}. This non-constructive
character is the price of considerable strength: choice, or one of its
equivalents, is what licenses such everyday principles as ``every vector space
has a basis'' and ``every surjection has a right inverse'', and the bulk of
modern mathematics is developed within ZFC.

\paragraph{Choice and its equivalents.}
The form of~(\ref{ax:choice}) above is due to Zermelo. Over the axioms of ZF it is
equivalent to the existence of a \emph{choice function} for every family of
non-empty sets, to Zorn's lemma~\cite{zorn1935remark}, and to the well-ordering
theorem; we use whichever formulation is convenient. The model of
\Cref{ch:zflean} validates choice, so the developments of this thesis are free to
appeal to it.

\subsection{Natural Numbers, Ordinals, and the Cumulative Hierarchy}
\label{subsec:set-theory-hierarchy}

The axiom of infinity~(\ref{ax:inf}) provides an inductive set, from which the
natural numbers are constructed.
Following von Neumann, each natural number is identified with the set of its
predecessors: $0 \triangleq \varnothing$ and $n + 1 \triangleq n \cup \{n\}$,
so that $n = \{0, 1, \dots, n - 1\}$ and $m < n \iff m \in n$.
The least inductive set\footnotemark---\ie the intersection of all inductive sets---is
the first infinite von Neumann ordinal $\regnotation{\omega} = \{0, 1, 2, \dots\}$,
also known as the set \(\mathbb{N}\) of natural numbers.

\footnotetext{
  The term \emph{inductive set} has carried different, inequivalent meanings.
  We use it in the Zermelo--von Neumann sense: a set containing $\varnothing$ and
  closed under the successor $x \mapsto x \cup \{x\}$, from which the natural
  numbers arise as transitive sets~\cite{zermelo1908untersuchungen,vonNeumann1925axiomatisierung}.
  In order theory, by contrast, Bourbaki calls an ordered set inductive when every
  chain has an upper bound---the hypothesis of Zorn's lemma~\cite{bourbaki2007}.
}

\begin{example}[von Neumann naturals]
The first few naturals are
$0 = \varnothing$, $\;1 = \{0\} = \{\varnothing\}$, $\;2 = \{0, 1\} = \{\varnothing, \{\varnothing\}\}$, $\;3 = \{0, 1, 2\} = \{\varnothing, \{\varnothing\}, \{\varnothing, \{\varnothing\}\}\}$. Order is
membership: $2 < 3$ because $2 \in 3$, and $2 \subseteq 3$ since every
predecessor of $2$ is a predecessor of $3$.
\end{example}

More generally, the \emph{ordinals} $\regnotation{\Ord}$ are the transitive sets
well-ordered by~$\in$, extending the naturals into the transfinite. Cardinality
measures the size of a set: two sets have the same \emph{cardinality} when their
elements can be put in one-to-one correspondence, and a \emph{cardinal} is an
ordinal measuring such a size (the least ordinal of its cardinality; the infinite
cardinals are the alephs $\aleph_0 = \omega, \aleph_1, \aleph_2, \dots$).
The ordinals index a \emph{cumulative hierarchy} of set universes, which we define after
Werner~\cite{werner1997sets} by a single transfinite recursion,
\[
  \regnotation{\cumhier}_\alpha \triangleq \bigcup_{\beta < \alpha} \pow(\cumhier_\beta).
\]
The empty union of the first stage gives $\cumhier_0 = \varnothing$; the recursion unfolds to
$\cumhier_{\alpha + 1} = \pow(\cumhier_\alpha)$ at successors and to
$\cumhier_\lambda = \bigcup_{\beta < \lambda} \cumhier_\beta$ at \emph{limit
ordinals} $\lambda$---those with no immediate predecessor---so a single clause
subsumes the three cases. The class $\cumhier = \bigcup_{\alpha \in \Ord} \cumhier_\alpha$ is, by foundation, the
entire universe: every set $X$ occurs at some stage, the least $\alpha$ with
$X \in \cumhier_\alpha$ being its \emph{rank}\footnotemark. This well-founded, iterative picture is
what the Lean model of \Cref{ch:zflean} renders constructive, representing sets
as an inductively generated, $\in$-well-founded type quotiented by extensional
equality.

\footnotetext{
  Different conventions may be found in the literature. Some authors take the least
  \(\alpha\) such that \(X \subseteq \cumhier_\alpha\) as the rank of \(X\)~\cite{jech2003set,kunen2011set}.
}

\begin{remark}[Universes as models of ZFC]
A striking feature of this hierarchy is Werner's Lemma~4~\cite{werner1997sets}: if $\kappa$ is an
\emph{inaccessible cardinal}---one so large that $\cumhier_\kappa$ is closed
under power set and small unions---then $\cumhier_\kappa$ itself satisfies every
axiom of ZFC, a set that models set theory in its own right (a \emph{Grothendieck
universe}). This is the hinge of Werner's mutual encodings between the Calculus
of Inductive Constructions and ZFC, matching type-theoretic universes with
inaccessible cardinals. It is also what legitimizes
developing genuine ZFC set theory inside the universe hierarchy of a
type-theoretic proof assistant, as we do in \Cref{ch:zflean}---a circle of ideas
that directly shaped the present work.
\end{remark}

These considerations already land beyond the scope of this thesis, and are purely contextual.
The axioms \eqref{ax:ext}--\eqref{ax:choice} above are all that the present work requires of its foundation.
The B method, to which we now turn (\Cref{sec:b-method}), is built directly upon them:
its expressions denote sets, and its relations, functions, and sequences are
particular sets of pairs, assembled into a hierarchy of function spaces. We accordingly present the B
language and the proof obligations it generates against this set-theoretic
background.
